% ============================================================
%  ArcLoom SPPU ASIC Architecture Specification
%  Version 1.0 — April 2026
%  Confidential — Internal Use Only
% ============================================================
\documentclass[11pt, a4paper]{article}

\usepackage[margin=2.2cm]{geometry}
\usepackage{amsmath, amssymb}
\usepackage{booktabs}
\usepackage{array}
\usepackage{longtable}
\usepackage{xcolor}
\usepackage{hyperref}
\usepackage{titlesec}
\usepackage{enumitem}
\usepackage{fancyhdr}
\usepackage{graphicx}

\hypersetup{colorlinks=true, linkcolor=black, urlcolor=blue}

\definecolor{tfe-gold}{RGB}{180, 140, 60}
\definecolor{tfe-dark}{RGB}{30, 30, 30}
\definecolor{tfe-gray}{RGB}{100, 100, 100}

\pagestyle{fancy}
\fancyhf{}
\rhead{\textcolor{tfe-gray}{\small ArcLoom SPPU ASIC Architecture v1.0}}
\lhead{\textcolor{tfe-gray}{\small Confidential --- Internal Use Only}}
\rfoot{\textcolor{tfe-gray}{\small Page \thepage}}

\titleformat{\section}{\large\bfseries\color{tfe-dark}}{}{0em}{}[\titlerule]
\titleformat{\subsection}{\normalsize\bfseries\color{tfe-dark}}{}{0em}{}
\titleformat{\subsubsection}{\normalsize\itshape\color{tfe-dark}}{}{0em}{}

\setlength{\parindent}{0pt}
\setlength{\parskip}{6pt}

\begin{document}

\begin{center}
  {\LARGE \textbf{ArcLoom SPPU}}\\[6pt]
  {\large ASIC Architecture Specification}\\[4pt]
  {\large Ternary Structural Perception Processor}\\[8pt]
  {\color{tfe-gray}\small Version 1.0 \quad|\quad April 2026 \quad|\quad Confidential}\\[4pt]
  {\color{tfe-gray}\small Joseph Forrester \quad|\quad DSF-AI Architecture}
\end{center}

\vspace{1em}
\hrule
\vspace{1em}

% ============================================================
\section{Executive Summary}

The ArcLoom SPPU (Structural Perception Processing Unit) is a novel computing
architecture based on balanced ternary logic that performs structural
perception through combinational trit coupling. Unlike conventional
binary processors that execute sequential instructions, the SPPU
settles to decisions through physical wire propagation in a coupled
trit fabric. No clock is required in the decision path. Power
consumption between sensor events is zero.

This document specifies the ASIC architecture for a multi-tile SPPU
chip suitable for fabrication in modern CMOS processes (28nm--5nm).

\textbf{Key differentiators:}
\begin{itemize}[nosep]
  \item Zero dynamic power between evaluations (no clock tree)
  \item Nanosecond decision latency (combinational settling)
  \item Fully deterministic (identical inputs $\rightarrow$ identical outputs)
  \item No machine learning, no probability, no heuristics
  \item Domain-agnostic (coupling weights define the application)
  \item Proven on FPGA: 19,683 arithmetic operations with zero errors,
        live sensor-to-decision demonstrated on Zynq-7020
\end{itemize}

% ============================================================
\section{Architecture Overview}

\subsection{System Hierarchy}

\begin{center}
\begin{tabular}{@{}llll@{}}
\toprule
\textbf{Level} & \textbf{Component} & \textbf{Count per chip} & \textbf{Function} \\
\midrule
Die       & ArcLoom SoC       & 1         & Complete system \\
Cluster   & Perception Cluster & 16--64   & Group of coupled SPPUs \\
Tile      & SPPU Tile          & 1,024+   & Single perception unit \\
Sublevel  & MathLoom ALU       & per tile  & Ternary arithmetic \\
Sublevel  & Krimelack Block    & per cluster & Structural memory \\
Sublevel  & BSIL Port          & perimeter & Sensor interface \\
Sublevel  & L6 Aggregator      & per cluster & Inevitability detection \\
\bottomrule
\end{tabular}
\end{center}

\subsection{Die Floorplan}

\begin{verbatim}
+----------------------------------------------------------+
|                    BSIL Sensor Ring                        |
|  +----------------------------------------------------+  |
|  |  Cluster 0  |  Cluster 1  |  Cluster 2  | Cluster 3|  |
|  |  [16 SPPU]  |  [16 SPPU]  |  [16 SPPU]  | [16 SPPU]|  |
|  |  Krimelack   |  Krimelack   |  Krimelack  | Krimelack|  |
|  |  L6 Aggr     |  L6 Aggr     |  L6 Aggr    | L6 Aggr  |  |
|  +----------------------------------------------------+  |
|  |  Cluster 4  |  Cluster 5  |  Cluster 6  | Cluster 7|  |
|  |  [16 SPPU]  |  [16 SPPU]  |  [16 SPPU]  | [16 SPPU]|  |
|  |  Krimelack   |  Krimelack   |  Krimelack  | Krimelack|  |
|  |  L6 Aggr     |  L6 Aggr     |  L6 Aggr    | L6 Aggr  |  |
|  +----------------------------------------------------+  |
|  |         MathLoom Array   |   Config/Debug           |  |
|  |         [shared ALUs]    |   [JTAG, host I/F]       |  |
|  +----------------------------------------------------+  |
|                    BSIL Sensor Ring                        |
+----------------------------------------------------------+
\end{verbatim}

% ============================================================
\section{SPPU Tile}

\subsection{Architecture}

Each SPPU tile is a self-contained combinational perception unit.
It has no clock input. The decision appears on the output wires
within one propagation delay of the inputs changing.

\textbf{Trit encoding:} 2 wires per trit.
\begin{itemize}[nosep]
  \item \texttt{00} = 0 (null / undecided)
  \item \texttt{01} = +1 (positive / excite)
  \item \texttt{10} = -1 (negative / inhibit)
  \item \texttt{11} = invalid (never produced)
\end{itemize}

\subsection{Strand Configuration}

8 strands $\times$ 3 trits = 24 trits = 48 wires per SPPU.

\begin{center}
\begin{tabular}{@{}llll@{}}
\toprule
\textbf{Strand} & \textbf{Trits} & \textbf{Type} & \textbf{Source} \\
\midrule
Distance      & 3 & Input   & BSIL (distance sensor) \\
Direction     & 3 & Input   & BSIL (delta of distance) \\
Acceleration  & 3 & Input   & BSIL (second derivative) \\
Sensor A      & 3 & Input   & BSIL (domain-specific) \\
Sensor B      & 3 & Input   & BSIL (domain-specific) \\
Context       & 3 & Settling & Coupled from inputs \\
Momentum      & 3 & Settling & Coupled from inputs \\
Decision      & 3 & Output  & Coupled from all strands \\
\bottomrule
\end{tabular}
\end{center}

\subsection{Coupling Network}

The coupling network is a feedforward combinational circuit with
two levels:

\textbf{Level 1 (parallel):}
Context and Momentum strands each compute their 3 trits from the
15 input trits (5 input strands $\times$ 3 trits) through weighted
summation with dead-zone thresholding.

\textbf{Level 2:}
Decision strand computes its 3 trits from all 21 trits (5 input +
2 settling strands) through weighted summation with dead-zone
thresholding.

\textbf{Local field computation} for each output trit:
\begin{equation}
  h_j = \sum_{i} J_{ij} \cdot t_i + h_j^{ext}
\end{equation}
where $J_{ij}$ is the signed 8-bit coupling weight, $t_i \in \{-1, 0, +1\}$
is the input trit value, and $h_j^{ext}$ is the external field bias.

\textbf{Trit activation:}
\begin{equation}
  t_j = \begin{cases}
    +1 & \text{if } h_j > \tau \\
    -1 & \text{if } h_j < -\tau \\
    0  & \text{if } |h_j| \leq \tau
  \end{cases}
\end{equation}
where $\tau$ is the dead-zone threshold. The null state (0) is a
first-class output — it represents structural uncertainty, not
a failure to decide.

\subsection{Gate Count Estimate}

\begin{center}
\begin{tabular}{@{}lr@{}}
\toprule
\textbf{Component} & \textbf{Equivalent Gates} \\
\midrule
Level 1: 6 local field units (15 inputs each) & 1,200 \\
Level 2: 3 local field units (21 inputs each) & 900 \\
Dead-zone comparators (9 total) & 100 \\
Input/output buffers & 50 \\
\midrule
\textbf{Total per SPPU tile} & \textbf{$\sim$2,250} \\
\bottomrule
\end{tabular}
\end{center}

\subsection{Power Analysis}

\textbf{Static power:} Leakage only. In 7nm FinFET, approximately
0.1 $\mu$W per SPPU tile.

\textbf{Dynamic power:} Zero between sensor events. When inputs
change, a single wave of switching propagates through the coupling
network ($\sim$50 gate transitions). At 0.7V supply:
\begin{equation}
  P_{event} \approx 50 \times C_{gate} \times V_{dd}^2 \approx 50 \times 0.5\text{fF} \times 0.49\text{V}^2 \approx 12\text{ fJ per event}
\end{equation}

At 1,000 sensor events per second: $P_{dynamic} \approx 12$ pW per tile.

\textbf{Comparison:} A single ARM Cortex-M0 core running inference at
1,000 Hz draws approximately 10 mW. A single SPPU tile performing
the same perception task draws 12 pW. That is a factor of
$\sim$10$^{9}$ (one billion) lower power.

\subsection{Timing}

Propagation delay through the coupling network:
\begin{itemize}[nosep]
  \item Level 1 (15-input weighted sum + threshold): $\sim$2 ns in 7nm
  \item Level 2 (21-input weighted sum + threshold): $\sim$3 ns in 7nm
  \item Total input-to-decision: $\sim$5 ns
\end{itemize}

No clock uncertainty. No pipeline latency. No cache misses.
Deterministic 5 ns from sensor change to decision output.

% ============================================================
\section{MathLoom Arithmetic Unit}

\subsection{Architecture}

The MathLoom provides balanced ternary arithmetic for BSIL
encoding, threshold computation, and general-purpose calculation.
It is combinational (no clock) and shared across a cluster of
SPPU tiles.

\textbf{Operations:}
\begin{center}
\begin{tabular}{@{}lll@{}}
\toprule
\textbf{Operation} & \textbf{Method} & \textbf{Verified} \\
\midrule
Addition (N-digit) & Balanced ternary ripple-carry & 6,561/6,561 on FPGA \\
Multiplication (N$\times$N) & Shift-and-add with trit negation & 6,561/6,561 on FPGA \\
Comparison & Subtraction + MSB scan & 6,561/6,561 on FPGA \\
Negation & Bit-swap per trit & Trivial (combinational) \\
Absolute value & Conditional negation on MSB & Verified in simulation \\
\bottomrule
\end{tabular}
\end{center}

\subsection{Balanced Ternary Full Adder}

Core digit constraint:
\begin{equation}
  a_i + b_i + c_i = s_i + 3 \cdot c_{i+1}
\end{equation}

Deterministic tie-break rule (normative):
\begin{enumerate}[nosep]
  \item Minimize $|s_i|$ (prefer 0 over $\pm$1)
  \item Minimize $|c_{i+1}|$
\end{enumerate}

Truth table (all 7 possible input sums):
\begin{center}
\begin{tabular}{@{}rrr@{}}
\toprule
\textbf{Total} & \textbf{Sum} & \textbf{Carry} \\
\midrule
$-3$ & $0$ & $-1$ \\
$-2$ & $+1$ & $-1$ \\
$-1$ & $-1$ & $0$ \\
$0$  & $0$  & $0$ \\
$+1$ & $+1$ & $0$ \\
$+2$ & $-1$ & $+1$ \\
$+3$ & $0$  & $+1$ \\
\bottomrule
\end{tabular}
\end{center}

\subsection{Multiplication}

4-digit $\times$ 4-digit balanced ternary multiplication via
shift-and-add. For each digit of operand B:
\begin{itemize}[nosep]
  \item If digit $= +1$: add A shifted to position
  \item If digit $= -1$: subtract A shifted to position (add negation)
  \item If digit $= 0$: skip (zero contribution)
\end{itemize}

Result: 8-digit product. Four partial products summed by three
8-digit adders. Entirely combinational.

\subsection{Key Property: Truncation = Rounding}

In balanced ternary, truncating a number to N digits automatically
rounds to the nearest representable value. Binary requires explicit
rounding logic. This property makes balanced ternary inherently
superior for signal processing and sensor data applications where
precision must be traded for efficiency.

\subsection{Folding Division}

MathLoom implements division using a novel method called
\textbf{folding division}. Given numerator $A$ and denominator $B$,
the method iteratively subtracts $|B|$ from $|A|$ until the residual
is less than $|B|$. The iteration count is the quotient magnitude.
The residual is the remainder. The sign is determined by the operand
signs.

\begin{equation}
  A = q \cdot B + r, \quad 0 \leq r < |B|
\end{equation}

\textbf{Native rounding:} When the remainder $r \geq |B|/2$, the
quotient is incremented by one. This produces nearest-integer
rounding as a natural consequence of the balanced ternary
representation, with no additional rounding logic or mode bits
required. This is distinct from all conventional division methods
(restoring, non-restoring, SRT, Newton-Raphson, Goldschmidt),
which require separate rounding stages.

\textbf{Fixed-point extension:} For fractional results, the numerator
is pre-scaled by $3^N$ before folding, where $N$ is the desired number
of fractional trit digits. This produces $N$ digits of fractional
precision. Note that $1/3$ is \textbf{exact} in balanced ternary
(unlike binary, which cannot represent $1/3$ finitely).

\textbf{Computational completeness:} Addition, multiplication, and
folding division together form a complete set of arithmetic
primitives. Any numerical function can be computed from these three
operations:
\begin{itemize}[nosep]
  \item Fixed-point fractions (scale and fold)
  \item Complex number arithmetic (paired real/imaginary folding)
  \item Matrix operations (element-wise folding)
  \item Polynomial arithmetic (synthetic folding)
  \item Square roots (Newton iteration: $x_{n+1} = (x_n + A/x_n) / 2$)
  \item Trigonometric functions (Taylor series or CORDIC)
\end{itemize}

\textbf{Verification:}
\begin{center}
\begin{tabular}{@{}llr@{}}
\toprule
\textbf{Category} & \textbf{Tests} & \textbf{Failures} \\
\midrule
Integer division & 6,480 & 0 \\
Complex number division & 14,520 & 0 \\
2$\times$2 matrix inversion & 6,016 & 0 \\
Polynomial division & 14,406 & 0 \\
Edge cases & 8 & 0 \\
\midrule
\textbf{Total} & \textbf{41,430} & \textbf{0} \\
\bottomrule
\end{tabular}
\end{center}

\subsection{Gate Count}

\begin{center}
\begin{tabular}{@{}lr@{}}
\toprule
\textbf{Component} & \textbf{Equivalent Gates} \\
\midrule
4-digit adder (4 full adders) & 60 \\
4$\times$4 multiplier (4 negators + 3 8-digit adders) & 400 \\
4-digit comparator (subtractor + MSB scan) & 80 \\
\midrule
\textbf{Total per MathLoom ALU} & \textbf{$\sim$540} \\
\bottomrule
\end{tabular}
\end{center}

% ============================================================
\section{BSIL — Binary Sensor Ingestion Layer}

\subsection{Function}

BSIL converts binary sensor data (analog voltages, digital buses,
serial protocols) into balanced ternary trit strands suitable for
the SPPU coupling fabric.

\subsection{Encoding Method}

For analog sensors (ADC input):
\begin{equation}
  t_i = \begin{cases}
    +1 & \text{if } v > \theta_{hi,i} \\
    -1 & \text{if } v < \theta_{lo,i} \\
    0  & \text{if } \theta_{lo,i} \leq v \leq \theta_{hi,i}
  \end{cases}
\end{equation}

Each input strand uses three threshold pairs, creating a multi-scale
encoding of the sensor value. The direction strand encodes the first
derivative (delta between consecutive samples). The acceleration
strand encodes the second derivative.

\subsection{ASIC Implementation}

\begin{itemize}[nosep]
  \item On-chip SAR ADC (12-bit, $<$1 $\mu$W) per BSIL port
  \item Configurable threshold registers (loaded at initialization)
  \item Delta and acceleration computation (one subtractor + registers)
  \item Clocked at sensor sample rate only (not system clock)
\end{itemize}

BSIL is the \textbf{only clocked component} in the perception path.
The clock runs at the sensor sample rate (typically 10--10,000 Hz),
not at GHz frequencies. Between samples, the clock stops and BSIL
draws zero dynamic power.

% ============================================================
\section{Krimelack — Structural Memory}

\subsection{Function}

Krimelack stores committed structural motifs (settled SPPU states
that met quality criteria) and provides cue-based recall for
pattern recognition across time.

\subsection{Architecture}

\begin{itemize}[nosep]
  \item Content-addressable memory (CAM) for parallel pattern matching
  \item 32--256 motif entries per cluster
  \item Each entry: 48 bits (24 trits) + 8-bit hash for duplicate detection
  \item Commit gating: motifs stored only when uncertainty $U^* \leq U_{settle}$
        and resonance $R \geq R_{commit}$ (prevents storing noise)
  \item Recall: parallel comparison against all entries, returns best match
        and match score
\end{itemize}

\subsection{Commit Eligibility (Normative)}

A motif is committed to Krimelack only when:
\begin{enumerate}[nosep]
  \item Settled window holds ($U^* \leq 0.35$ for $K_s$ consecutive gates)
  \item Resonance exceeds threshold ($R \geq 0.72$)
  \item SafeMode is not active
  \item No duplicate exists (idempotency)
\end{enumerate}

\subsection{Gate Count}

\begin{center}
\begin{tabular}{@{}lr@{}}
\toprule
\textbf{Component} & \textbf{Equivalent Gates} \\
\midrule
CAM array (64 entries $\times$ 48 bits) & 4,000 \\
Hash comparison logic & 200 \\
Commit control + counter & 100 \\
Best-match scorer & 300 \\
\midrule
\textbf{Total per Krimelack block} & \textbf{$\sim$4,600} \\
\bottomrule
\end{tabular}
\end{center}

% ============================================================
\section{L6 — Topological Constraint Layer}

\subsection{Function}

L6 monitors the SPPU loom state and detects when dimensional
exhaustion has collapsed the manifold of possible states to a
singularity. When L6 fires a Structural Lock (SL-1), the decision
is not probable — it is geometrically inevitable.

\subsection{Implementation}

Purely combinational. Counts the number of non-null trits
(collapsed dimensions) in the loom state.

\textbf{Structural Lock criterion:}
\begin{equation}
  n_{effective} < \frac{n_{start}}{e}
\end{equation}

For 24 trits: $n_{start} = 24$, knee at $24/e \approx 8.83$.
SL-1 fires when 16 or more trits are non-null ($n_{effective} < 9$).

The threshold is \textbf{not a design parameter}. It is a mathematical
consequence of how volume behaves in $n$-dimensional space (the
gamma function inflection point). See L6-TCL Specification v1.0.

\subsection{Gate Count}

\begin{center}
\begin{tabular}{@{}lr@{}}
\toprule
\textbf{Component} & \textbf{Equivalent Gates} \\
\midrule
Trit null detector (24 comparators) & 48 \\
Population counter (5-bit) & 60 \\
Threshold comparator & 10 \\
Disruption gate & 5 \\
\midrule
\textbf{Total per L6 unit} & \textbf{$\sim$123} \\
\bottomrule
\end{tabular}
\end{center}

% ============================================================
\section{Perception Cluster}

\subsection{Architecture}

A Perception Cluster groups 16 SPPU tiles with shared Krimelack
memory and an L6 aggregator. SPPUs within a cluster can couple
to their nearest neighbors through a local coupling mesh.

\textbf{Inter-SPPU coupling:} Each SPPU's decision strand can
feed into neighboring SPPUs as additional input strands, creating
a lattice of coupled perception units. This enables multi-scale
structural perception — individual SPPUs perceive local features,
clusters perceive patterns, the full chip perceives scenes.

\subsection{Cluster Gate Count}

\begin{center}
\begin{tabular}{@{}lrr@{}}
\toprule
\textbf{Component} & \textbf{Per unit} & \textbf{$\times$ Count} \\
\midrule
SPPU tiles & 2,250 & $\times$ 16 = 36,000 \\
MathLoom ALU (shared) & 540 & $\times$ 1 = 540 \\
Krimelack block & 4,600 & $\times$ 1 = 4,600 \\
L6 aggregator & 123 & $\times$ 1 = 123 \\
Coupling mesh & 500 & $\times$ 1 = 500 \\
\midrule
\textbf{Total per cluster} & & \textbf{$\sim$41,763} \\
\bottomrule
\end{tabular}
\end{center}

% ============================================================
\section{Full Chip Specification}

\subsection{Configuration: ArcLoom-256}

\begin{center}
\begin{tabular}{@{}lr@{}}
\toprule
\textbf{Parameter} & \textbf{Value} \\
\midrule
SPPU tiles & 256 (16 clusters $\times$ 16 tiles) \\
MathLoom ALUs & 16 (one per cluster) \\
Krimelack blocks & 16 (one per cluster, 64 motifs each) \\
L6 aggregators & 16 + 1 global \\
BSIL ports & 32 (perimeter) \\
Total equivalent gates & $\sim$700,000 \\
\midrule
\textbf{Process target} & 28nm CMOS \\
\textbf{Die area estimate} & $\sim$1.5 mm$^2$ \\
\textbf{Package} & QFN-48 or BGA \\
\textbf{Supply voltage} & 0.9V (core) / 1.8V (I/O) \\
\textbf{Static power} & $<$ 50 $\mu$W \\
\textbf{Dynamic power (1kHz sensing)} & $<$ 5 $\mu$W \\
\textbf{Decision latency} & $<$ 10 ns \\
\textbf{Sensor-to-decision (with BSIL)} & $<$ 1 $\mu$s \\
\bottomrule
\end{tabular}
\end{center}

\subsection{Configuration: ArcLoom-4K}

\begin{center}
\begin{tabular}{@{}lr@{}}
\toprule
\textbf{Parameter} & \textbf{Value} \\
\midrule
SPPU tiles & 4,096 (256 clusters $\times$ 16 tiles) \\
Total equivalent gates & $\sim$11M \\
\textbf{Process target} & 7nm FinFET \\
\textbf{Die area estimate} & $\sim$3 mm$^2$ \\
\textbf{Static power} & $<$ 500 $\mu$W \\
\textbf{Dynamic power (1kHz sensing)} & $<$ 50 $\mu$W \\
\bottomrule
\end{tabular}
\end{center}

\subsection{Configuration: ArcLoom-1M (Vision)}

\begin{center}
\begin{tabular}{@{}lr@{}}
\toprule
\textbf{Parameter} & \textbf{Value} \\
\midrule
SPPU tiles & 1,048,576 (1M) \\
Perception clusters & 65,536 \\
\textbf{Process target} & 5nm / 3nm \\
\textbf{Die area estimate} & $\sim$50 mm$^2$ \\
\textbf{Application} & Multi-modal sensory cognition \\
\bottomrule
\end{tabular}
\end{center}

% ============================================================
\section{Comparison with Existing Architectures}

\begin{center}
\begin{tabular}{@{}lllll@{}}
\toprule
\textbf{Property} & \textbf{ARM Cortex-M0} & \textbf{Neural Accel.} & \textbf{SPPU (this work)} \\
\midrule
Computing model & Sequential & Matrix multiply & Combinational settling \\
Clock required & Yes (MHz--GHz) & Yes (MHz) & No (decision path) \\
Power (idle) & $\mu$W--mW & $\mu$W--mW & pW--nW \\
Power (active) & mW & mW & pW--$\mu$W \\
Decision latency & $\mu$s--ms & $\mu$s & ns \\
Deterministic & No (interrupts) & No (floating point) & Yes (combinational) \\
ML required & N/A & Yes & No \\
Uncertainty repr. & Computed & Probability & Native (null trit) \\
Rounding & Explicit & Explicit & Free (truncation) \\
\bottomrule
\end{tabular}
\end{center}

% ============================================================
\section{Intellectual Property Strategy}

\subsection{Patentable (Non-Trade-Secret)}

\begin{enumerate}[nosep]
  \item Combinational trit coupling fabric with dead-zone thresholding
  \item BSIL binary-to-ternary sensor encoding with multi-scale thresholds
  \item L6 topological constraint detection via Euler knee counting
  \item 8-strand SPPU topology with input/settling/decision strand hierarchy
  \item Balanced ternary shift-and-add multiplication with trit negation
  \item Krimelack content-addressable structural memory with commit gating
  \item Inter-SPPU coupling mesh for multi-scale perception
  \item Folding division method for balanced ternary arithmetic with
        native free rounding via truncation
\end{enumerate}

\subsection{Trade Secret (Protected)}

\begin{enumerate}[nosep]
  \item Coupling weight derivation methodology
  \item UF pipeline (L0--L4) internal formulas and thresholds
  \item DSF tuple composition and feedback mechanism
  \item Coherence field physics (Not-Math / UFCP)
  \item Kernel optimization procedures
\end{enumerate}

% ============================================================
\section{Validation Status}

\begin{center}
\begin{tabular}{@{}llll@{}}
\toprule
\textbf{Claim} & \textbf{Status} & \textbf{Evidence} \\
\midrule
Ternary addition correct & \textbf{PROVEN} & 6,561/6,561 on FPGA \\
Ternary multiplication correct & \textbf{PROVEN} & 6,561/6,561 on FPGA \\
Ternary comparison correct & \textbf{PROVEN} & 6,561/6,561 on FPGA \\
Combinational loom settling & \textbf{PROVEN} & Live sensor on FPGA \\
L6 structural lock detection & \textbf{PROVEN} & Demonstrated on FPGA \\
8-strand SPPU decisions & \textbf{PROVEN} & Live sensor on FPGA \\
BSIL sensor-to-trit encoding & \textbf{PROVEN} & Distance/direction/accel on FPGA \\
Krimelack commit gating & \textbf{PROVEN} & Motif storage verified on FPGA \\
Folding division correctness & \textbf{PROVEN} & 41,430/41,430 in simulation \\
Computational completeness & \textbf{PROVEN} & Complex, matrix, polynomial, trig \\
Direct HW sensor path (no CPU) & \textbf{PENDING} & Awaiting Pmod AD5 \\
Motor actuation from decisions & \textbf{PROVEN} & TB6612FNG, all 4 directions verified \\
Camera BSIL encoding & \textbf{PENDING} & Camera not connected \\
Domain agnosticism & \textbf{PENDING} & Single sensor type tested \\
Multi-SPPU coupling & \textbf{UNPROVEN} & Architecture defined, not built \\
\bottomrule
\end{tabular}
\end{center}

% ============================================================
\section{Target Applications}

\subsection{Tier 1: Immediate (FPGA prototype)}
\begin{itemize}[nosep]
  \item Autonomous micro-robotics (sensor-to-motor, zero CPU)
  \item Industrial edge monitoring (vibration, temperature, pressure)
\end{itemize}

\subsection{Tier 2: Near-term (ArcLoom-256 ASIC)}
\begin{itemize}[nosep]
  \item Implantable medical sensors (cardiac, glucose, neural)
  \item Environmental sensor networks (decade battery life)
  \item Safety-critical systems (deterministic, provable behavior)
\end{itemize}

\subsection{Tier 3: Long-term (ArcLoom-4K and beyond)}
\begin{itemize}[nosep]
  \item Multi-modal sensory perception (vision + audio + tactile)
  \item Nanoscale controllers (molecular machines, catalysts)
  \item Autonomous systems (drone swarms, underwater vehicles)
\end{itemize}

\subsection{Tier 4: Vision (ArcLoom-1M+)}
\begin{itemize}[nosep]
  \item Artificial sensory cognition
  \item Brain-scale structural perception
  \item DSF-AI universal cognition engine
\end{itemize}

% ============================================================
\section{Development Roadmap}

\begin{center}
\begin{tabular}{@{}lll@{}}
\toprule
\textbf{Phase} & \textbf{Milestone} & \textbf{Deliverable} \\
\midrule
Phase 0 (current) & FPGA proof of concept & PYNQ-Z2 demo \\
Phase 1 & Direct HW sensor path & Pmod AD5 integration \\
Phase 2 & Multi-sensor demo & Camera + distance + motor \\
Phase 3 & Patent filing & Architecture claims \\
Phase 4 & ASIC design (ArcLoom-256) & RTL + synthesis scripts \\
Phase 5 & Tape-out & Shuttle run (e.g., Efabless) \\
Phase 6 & Silicon validation & Chip-on-board demo \\
\bottomrule
\end{tabular}
\end{center}

% ============================================================
\section{Version History}

\begin{center}
\begin{tabular}{@{}lll@{}}
\toprule
\textbf{Version} & \textbf{Date} & \textbf{Description} \\
\midrule
1.0 & 2026-04-19 & Initial specification. FPGA-validated architecture. \\
\bottomrule
\end{tabular}
\end{center}

\vspace{2em}
\hrule
\vspace{0.5em}
{\color{tfe-gray}\small This document is confidential and for internal use only.
The ArcLoom SPPU architecture is the intellectual property of
Joseph Forrester / DSF-AI. The coupling weight derivation methodology,
UF pipeline internals, and coherence field physics are trade secrets
and are not disclosed in this document.}

\end{document}
